\documentclass[11pt]{article}

\usepackage[margin=1in]{geometry}
\usepackage{amsmath}
\usepackage{booktabs}
\usepackage{array}
\usepackage{enumitem}
\usepackage{xcolor}
\usepackage{listings}

\lstset{
    basicstyle=\ttfamily\small,
    breaklines=true,
    frame=single,
    columns=fullflexible
}

\title{Q-Learning Based Optimization for 3.5D NoC}
\author{}
\date{}

\begin{document}

\maketitle

\section{Goal}

The goal of the Q-learning optimizer is to find a good 3.5D NoC placement.
A placement contains:

\begin{itemize}[noitemsep]
    \item core-to-router mapping,
    \item TSV placement,
    \item TSV assignment,
    \item chiplet attach router.
\end{itemize}

The optimizer starts from random valid placements and repeatedly applies small
changes. Q-learning learns which type of change is useful in each search
situation.

\section{Chromosome Representation}

For a 3D chiplet, one gene is:

\[
[\text{core\_to\_router},\ \text{tsv\_placement},\ \text{tsv\_assignment},\ \text{base\_attach\_router}]
\]

For a 2.5D chiplet, one gene is:

\[
[\text{core\_to\_router},\ \text{base\_attach\_router}]
\]

The complete chromosome is a list of genes, one per chiplet.

\section{Actions}

At each iteration, Q-learning chooses one action from the following set:

\[
A =
\{
\text{swap\_mapping},
\text{reverse\_mapping},
\text{flip\_tsv},
\text{move\_attach}
\}
\]

\begin{table}[h]
\centering
\begin{tabular}{>{\ttfamily}l p{0.65\linewidth}}
\toprule
Action & Meaning \\
\midrule
swap\_mapping & Swap two cores between two router positions. \\
reverse\_mapping & Reverse a continuous block of the core-to-router mapping. \\
flip\_tsv & Add or remove one TSV position by flipping one TSV bit. \\
move\_attach & Change the router through which a chiplet connects to the base/interposer. \\
\bottomrule
\end{tabular}
\caption{Q-learning actions used by the optimizer.}
\end{table}

\subsection{Example: swap\_mapping}

Before:

\[
\begin{array}{c|cccc}
\text{Router} & 0 & 1 & 2 & 3 \\
\hline
\text{Core} & A & B & C & D
\end{array}
\]

Swap router 1 and router 3:

\[
\begin{array}{c|cccc}
\text{Router} & 0 & 1 & 2 & 3 \\
\hline
\text{Core} & A & D & C & B
\end{array}
\]

\subsection{Example: flip\_tsv}

Suppose TSV placement is:

\[
[1,\ 0,\ 1,\ 0]
\]

If position 1 is flipped:

\[
[1,\ 1,\ 1,\ 0]
\]

If position 2 is flipped:

\[
[1,\ 0,\ 0,\ 0]
\]

After every action, the chromosome is repaired to satisfy mapping and TSV
constraints.

\section{State}

The Q-learning state tells the agent the current search situation. The state is:

\[
s = (\text{fitness\_bucket},\ \text{stall\_bucket})
\]

It is computed from two values:

\[
\text{relative\_gap}
=
\frac{\text{current\_fitness} - \text{best\_fitness}}
{\text{best\_fitness}}
\]

and

\[
\text{stall\_count}
=
\text{number of iterations without improvement}
\]

\subsection{Fitness Bucket}

\[
\text{fitness\_bucket} =
\begin{cases}
0, & \text{relative\_gap} < 0.01 \\
1, & \text{relative\_gap} < 0.05 \\
2, & \text{relative\_gap} < 0.15 \\
3, & \text{otherwise}
\end{cases}
\]

\subsection{Stall Bucket}

\[
\text{stall\_bucket} =
\begin{cases}
0, & \text{stall\_count} < 10 \\
1, & \text{stall\_count} < 40 \\
2, & \text{stall\_count} < 100 \\
3, & \text{otherwise}
\end{cases}
\]

Example:

\[
s = (0,0)
\]

means the current solution is very close to the best solution and the search is
not stuck.

\[
s = (3,3)
\]

means the current solution is far from the best solution and the search is badly
stuck.

\section{Q-Table}

The Q-table stores one score for each state-action pair:

\[
Q(s,a)
\]

Example:

\begin{table}[h]
\centering
\begin{tabular}{c c c c c}
\toprule
State & swap\_mapping & reverse\_mapping & flip\_tsv & move\_attach \\
\midrule
$(0,0)$ & 0.80 & 0.20 & -0.10 & 0.05 \\
$(1,2)$ & 0.05 & 0.10 & 0.85 & 0.40 \\
$(3,3)$ & -0.10 & 0.25 & 0.60 & 0.75 \\
\bottomrule
\end{tabular}
\caption{Example Q-table.}
\end{table}

If the current state is $(1,2)$, the best action is \texttt{flip\_tsv}, because
it has the highest Q-value:

\[
Q((1,2), \text{flip\_tsv}) = 0.85
\]

\section{Action Selection}

The optimizer uses epsilon-greedy action selection.

With probability $\epsilon$, it explores:

\[
a = \text{random action}
\]

With probability $1-\epsilon$, it exploits:

\[
a =
\arg\max_{a'} Q(s,a')
\]

At the beginning, $\epsilon$ is high, so the optimizer explores more. Over time,
$\epsilon$ decays, so the optimizer uses learned actions more often.

\section{Reward}

After applying an action, the candidate fitness is computed.

\[
\text{improvement}
=
\text{current\_fitness} - \text{candidate\_fitness}
\]

The reward is:

\[
r =
\frac{\text{improvement}}
{|\text{current\_fitness}|}
\]

If the candidate becomes the best solution found so far, an extra reward is
added:

\[
r = r + 1
\]

If the candidate is worse or equal, a small penalty is applied:

\[
r = r - 0.02
\]

\section{Q-Value Update}

The Q-table is updated using the standard Q-learning equation:

\[
Q(s,a)
\leftarrow
Q(s,a)
 +
\alpha
\left[
r
 +
\gamma \max_{a'} Q(s',a')
 -
Q(s,a)
\right]
\]

where:

\begin{itemize}[noitemsep]
    \item $s$ is the current state,
    \item $a$ is the selected action,
    \item $r$ is the reward,
    \item $s'$ is the next state,
    \item $\alpha$ is the learning rate,
    \item $\gamma$ is the discount factor.
\end{itemize}

\section{Algorithm Flow}

\begin{enumerate}
    \item Build the 3.5D topology.
    \item Create the edge set from the input graph.
    \item Generate a random initial population.
    \item Pick the best initial chromosome as the current solution.
    \item Initialize the Q-table as empty.
    \item For each iteration:
    \begin{enumerate}
        \item Compute current state.
        \item Select an action using epsilon-greedy policy.
        \item Apply the selected action.
        \item Repair the chromosome.
        \item Compute candidate fitness.
        \item Compute reward.
        \item Update the Q-table.
        \item Update best solution if candidate is better.
        \item Decay epsilon.
    \end{enumerate}
    \item Write the best chromosome as a particle file.
\end{enumerate}

\section{Pseudocode}

\begin{lstlisting}
initialize Q-table as empty
generate initial chromosomes
current = best initial chromosome
best = current

for each iteration:
    state = get_state(current_fitness, best_fitness, stall_count)

    if random() < epsilon:
        action = random action
    else:
        action = action with highest Q(state, action)

    candidate = apply_action(current, action)
    candidate = repair(candidate)
    candidate_fitness = evaluate(candidate)

    reward = compute_reward(current_fitness, candidate_fitness)

    next_state = get_state(candidate_fitness, best_fitness, stall_count)

    Q(state, action) =
        Q(state, action)
        + alpha * (reward + gamma * max Q(next_state, *) - Q(state, action))

    if candidate is accepted:
        current = candidate

    if candidate is better than best:
        best = candidate
        stall_count = 0
    else:
        stall_count = stall_count + 1

    epsilon = max(epsilon_min, epsilon * epsilon_decay)

write best particle
\end{lstlisting}

\section{Summary}

The Q-learning optimizer learns which mutation action is useful in each search
situation. Instead of randomly applying changes throughout the run, it builds a
Q-table that remembers which actions improved fitness before.

\[
\text{Q-learning memory}
=
\text{which action works best in which state}
\]

\end{document}
